% ============================================================
% Modern Colorful 模板用法速查 — 展示所有可用环境和命令
% 编译方式: XeLaTeX
% ============================================================
\documentclass{ModernColorful}

\hypersetup{
  pdfauthor={作者名},
  pdftitle={论文标题 深度解读}
}

\title{\textbf{论文标题\\深度解读}}
\author{作者名}
\date{2026年3月17日}

\begin{document}

% === 标题页 ===
\begin{titlepage}
  \AddTitlePageLogo
  \thispagestyle{empty}
  \centering
  \vspace*{2cm}

  \begin{titlebox}
    \centering
    {\LARGE\bfseries\color{primaryblue} 论文标题}\\[10pt]
    {\Large\bfseries 论文深度解读}\\[15pt]
    {\normalsize 一句话副标题}
  \end{titlebox}

  \vspace{2cm}

  {\large\color{textgray}
  \textbf{作者}：作者名\\[5pt]
  \textbf{日期}：2026年3月17日
  }

  \vfill

  \begin{infobox}[论文信息]
    \centering
    \textbf{原文}：\href{https://arxiv.org/abs/XXXX.XXXXX}{arXiv:XXXX.XXXXX}\\
    \textbf{代码}：\href{https://github.com/user/repo}{github.com/user/repo}
  \end{infobox}
\end{titlepage}

\tableofcontents
\newpage

% ============================================================
\section{背景知识}
% ============================================================

正文段落示例，可使用\term{术语标记}来标注关键术语。

\subsection{子节标题}

\begin{warningbox}[注意事项]
  \begin{itemize}
    \item \hlblue{蓝色要点}：说明文字
    \item \hlorange{橙色要点}：说明文字
  \end{itemize}
\end{warningbox}

\subsubsection{子子节标题}

\begin{itemize}
  \item \hlblue{蓝色高亮}
  \item \hlpurple{紫色高亮}
  \item \hlteal{青色高亮}
  \item \hlorange{橙色高亮}
  \item \hlbox[bgpink]{背景高亮}
\end{itemize}

% ============================================================
\section{核心问题与创新思路}
% ============================================================

\begin{conceptbox}[概念盒子]
  用于解释背景概念或对比分析。
\end{conceptbox}

\begin{innovationbox}[创新点]
  用于展示核心创新贡献。
\end{innovationbox}

\begin{infobox}[信息盒子]
  \centering
  用于一般信息、实验配置等。
\end{infobox}

% ============================================================
\section{方法详解}
% ============================================================

重要公式使用公式盒子：

\begin{equationbox}
  \centering
  $E = mc^2$
\end{equationbox}

有序列表：
\begin{enumerate}
  \item \hlblue{步骤一}：说明
  \item \hlpurple{步骤二}：说明
\end{enumerate}

% ============================================================
\section{实验结果}
% ============================================================

\begin{center}
\begin{tabular}{>{\raggedright}p{3cm} >{\centering}p{3cm} >{\centering\arraybackslash}p{3cm}}
  \toprule
  \tableheadercolor
  \textbf{方法} & \textbf{指标A} & \textbf{指标B} \\
  \midrule
  基线方法 & 1.0 & 基线 \\
  \rowcolor{bggreen}
  \textbf{本文方法} & \textbf{2.5} & \textbf{最优} \\
  \bottomrule
\end{tabular}
\end{center}

代码/伪代码盒子：

\begin{codebox}[算法伪代码]
  1. Initialize model\\
  2. For each step: predict and verify\\
  3. Return result
\end{codebox}

图片（需先复制到报告目录）：

\begin{figure}[H]
  \centering
  % \includegraphics[width=0.7\textwidth]{figure.pdf}
  \caption{图片标题}
\end{figure}

% ============================================================
\section{总结}
% ============================================================

\begin{titlebox}
  \textbf{核心贡献}
  \begin{enumerate}
    \item 贡献一：\hlblue{关键词}
    \item 贡献二：\hlpurple{关键词}
  \end{enumerate}
\end{titlebox}

% ============================================================
\section{参考资料}
% ============================================================

\begin{enumerate}
  \item Author, A. (2026). Paper Title. Conference/Journal.
\end{enumerate}

\end{document}